\section{Feature Prototype and Evaluation}

This section presents a focused implementation of the explanation generator and a comparative evaluation of several local large language models (LLMs). The objective is to determine which model best aligns with the system’s requirements for explainability, resource efficiency, and accessibility. The evaluation is scoped to models runnable within an 8 GB VRAM constraint to support reproducibility and local deployment.

\subsection{Prototype Objective and Scope}

The prototype isolates the final layer of the GenAI Advisor architecture—the explanation generator. It takes as input a structured recommendation, typically generated by the strategy engine, and produces a textual rationale. This explanation layer is central to the system’s commitment to transparency and user trust. For the purposes of this prototype, upstream stages such as data collection, feature engineering, and scoring logic are simulated using manually specified inputs. This allows a controlled environment in which model outputs can be directly compared.

The task is framed as follows: given a set of financial signals and a recommendation (e.g., BUY, HOLD, SELL), can a language model explain the recommendation in a manner that is accurate, clear, and suitable for a non-technical retail investor? The success of the explanation layer depends on the model’s ability to infer causality or justification purely from the structured metrics it receives.

\subsection{Implementation Details}

The prototype is implemented in Python and interfaces with the \texttt{Ollama} runtime to serve quantised language models locally. Ollama is selected for its compatibility with consumer-grade hardware and its support for a growing list of well-maintained, open-source models. Using this framework, models can be served from the command line or accessed via HTTP API for more structured interaction.

All models tested are quantised 7B variants, chosen to balance performance with hardware constraints. Models include:

\begin{itemize}
    \item \texttt{mistral-7b} — A strong general-purpose instruction-tuned model known for its coherence and performance on reasoning tasks.
    \item \texttt{phi-7b} — Optimised by Microsoft for compactness and low hallucination rates.
    \item \texttt{gemma-7b} — Released by Google, this model offers efficient inference and strong factual alignment.
    \item \texttt{llama2-7b} — Meta’s widely adopted model, used here as a baseline.
    \item \texttt{deepseek-r1-7b} — A reasoning-oriented model with good handling of structured data.
    \item \texttt{vicuna-7b}, \texttt{codellama-7b}, \texttt{starling-lm-7b}, \texttt{neural-chat-7b}, and others — evaluated for comparative performance.
\end{itemize}

Each model is presented with a structured prompt template. No commentary or text is provided; the model must infer and generate all rationale. For example:

\begin{verbatim}
You are a financial assistant. Based on the following metrics:
Ticker: MSFT
Action: BUY
P/E Ratio: 32.1
RSI: 58
Beta: 1.1

Explain in simple terms to a layperson why this action is appropriate, using only the signals provided.
\end{verbatim}

Additional test prompts used during evaluation include:

\begin{verbatim}
You are a financial assistant. Based on the following metrics:
Ticker: NVDA
Action: HOLD
P/E Ratio: 75.6
RSI: 69
Beta: 1.8

Explain the rationale behind this HOLD recommendation to a retail investor.
\end{verbatim}

\begin{verbatim}
You are a financial assistant. Based on the following metrics:
Ticker: AMZN
Action: SELL
P/E Ratio: 92.4
RSI: 73
Beta: 1.4

Explain in accessible language why this stock may no longer be suitable to hold.
\end{verbatim}

Each of these prompts reflects a different combination of quantitative signals to test the model's consistency and interpretive ability across scenarios.

\subsection{Evaluation Methodology}

Evaluation focuses on three core criteria:

\begin{enumerate}
    \item \textbf{Factual Alignment}: Whether the explanation references and reflects the correct input signals (e.g., “high P/E” when P/E exceeds 80).
    \item \textbf{Interpretability}: Whether the rationale is readable, free from jargon, and suitable for a non-expert.
    \item \textbf{Stability}: Whether the same prompt produces consistent responses across repeated runs.
\end{enumerate}

Each model was tested on three prompts and scored manually. Ratings were assigned on a 1–5 scale per criterion and averaged across prompts. Qualitative impressions were recorded, focusing on how well each model translated technical indicators into investment-relevant language.

Responses were analysed for signal grounding, reasoning flow, and verbosity. Models that hallucinated commentary not supported by the input (e.g., "recent news coverage suggests...") were penalised.

\subsection{Results and Comparison}

The results are summarised in the following table:

\begin{center}
\begin{tabular}{lccc}
\textbf{Model} & \textbf{Factual Alignment} & \textbf{Interpretability} & \textbf{Stability} \\
\hline
\texttt{mistral-7b}        & 5/5 & 5/5 & 5/5 \\
\texttt{phi-7b}            & 5/5 & 4/5 & 5/5 \\
\texttt{gemma-7b}          & 4/5 & 4/5 & 4/5 \\
\texttt{llama2-7b}         & 3/5 & 3/5 & 4/5 \\
\texttt{vicuna-7b}         & 4/5 & 3/5 & 4/5 \\
\texttt{codellama-7b}      & 4/5 & 3/5 & 4/5 \\
\texttt{neural-chat-7b}    & 3/5 & 4/5 & 4/5 \\
\texttt{starling-lm-7b}    & 4/5 & 4/5 & 3/5 \\
\texttt{llama2-uncensored-7b} & 3/5 & 3/5 & 4/5 \\
\texttt{llava-7b}          & 4/5 & 3/5 & 4/5 \\
\texttt{deepseek-r1-7b}    & 5/5 & 3/5 & 4/5 \\
\end{tabular}
\end{center}

\texttt{mistral-7b} produced consistently high-quality explanations. It maintained clear structure, properly reflected all input values, and varied phrasing across runs. Explanations typically began with a signal summary and concluded with an action justification, closely resembling expert commentary. Its ability to adapt language depending on the action (e.g., cautious phrasing for HOLD or SELL) was particularly effective.

\texttt{phi-7b} showed similar strength in factual grounding but was often terse. It used precise language and made no unsupported claims, but the prose was more mechanical. \texttt{gemma-7b} was reliable and compact but lacked the fluency of \texttt{mistral-7b}.

Other models such as \texttt{vicuna-7b} and \texttt{starling-lm-7b} produced grammatically correct responses but were more generic. They occasionally repeated template phrasing or omitted one or more key signals (e.g., ignoring RSI or beta). These outputs would require post-processing or filtering before use in production.

\subsection{Discussion}

The comparison highlights how fine-tuned LLM behaviour affects trustworthiness in explainable AI. Models tuned for dialogue or code (\texttt{vicuna}, \texttt{codellama}) performed adequately but not optimally. Meanwhile, models like \texttt{mistral} and \texttt{phi}—both tuned for instruction following—demonstrated more appropriate structuring and interpretation.

The scoring rubric also reveals trade-offs. For example, \texttt{deepseek} was extremely accurate but slightly verbose and rigid in tone, which might reduce usability for lay audiences. \texttt{neural-chat} was fluent but tended toward overconfident or speculative claims.

Stability was assessed via repeated prompt submission. \texttt{mistral}, \texttt{phi}, and \texttt{gemma} retained output meaning while slightly varying sentence structure. This behaviour is desirable: the explanation remains fresh while still grounded.

\subsection{Conclusion}

This evaluation confirms that quantised, open-source 7B LLMs are capable of generating clear and accurate explanations from structured financial inputs. Among the models tested, \texttt{mistral-7b} is selected as the default explanation engine. It strikes a compelling balance between fluency, precision, and alignment with non-expert communication needs.

Its performance justifies its use as a core component of the GenAI Advisor system and supports the overall goal of delivering transparent, explainable financial recommendations. The modular design of the prototype allows future integration with upstream pipelines, where generated explanations will be contextualised within retrieved evidence and user-specific preferences.

This work illustrates a practical methodology for evaluating LLMs in structured domains and contributes to the broader discussion on scalable explainability for AI-assisted financial systems.
